\subsection{Magnetischer Kreis 1}
\Aufgabe{
    Berechnen Sie die die Induktivität $L$ der Ringkernspule in Abbildung \ref{AbbAufgMagnKreis1} mit rechteckförmigen Kernquerschnitt

	\f{width=\textwidth}{width=0.5\textwidth}{Übung_Ringspule_MagFeld.png}{Ringspule zu Aufgabe \label{AbbAufgMagnKreis1}}

	Die Spule hat die folgenden Daten:
	\begin{align*}
		N &= 250\\
		h &= 2\,\mathrm{cm}\\
		r_1 &= 3\,\mathrm{cm}\\
		r_2 &= 5\,\mathrm{cm}\\
		\mu_r &= 500\\
		\mu_0 &= 4\pi\cdot 10^{-7}\,\frac{\mathrm{Vs}}{\mathrm{Am}}
	\end{align*}

	\begin{enumerate}
		\item[\bf a)]
		mit Hilfe des magnetischen Widerstands $R_m$.
		\item[\bf b)]
		mit Hilfe der magnetischen Flussdichte $B$ und Feldstärke $H$.
		\item[\bf c)]
		Wie erklären Sie sich die abweichenden Ergebnisse aus Aufgabenteil a) und b)?
	\end{enumerate}
}

\Loesung{
	\begin{itemize}

		\item[a)]
		Magnetischer Widerstand:
		\begin{align*}
		R_{m} &= \frac{l_{Fe}}{\mu_r \cdot \mu_0 \cdot A_{Fe}} \\
		R_{m} &= \frac{2\pi \cdot (\frac{r_2 + r_1}{2})}{\mu_r \cdot \mu_0 \cdot h \cdot (r_2 - r_1)}
		\end{align*}

		Annahme: Im Kernquerschnitt $A$ verlaufe der magnetische Fluss $\Phi$ und das Magnetfeld habe überall die gleiche Flussdichte $B$.\\
		Mittlere Feldlinienlänge:
		\begin{equation*}
		\ell_{Fe} = 2\pi \cdot \frac{r_2 + r_1}{2}
		\end{equation*}

		Kernquerschnitt:
		\begin{equation*}
		A_{Fe} = h \cdot (r_2 - r_1)
		\end{equation*}

		Induktivität:
		\begin{align*}
		L_1 &= \frac{N^2}{R_{m}} = N^2 \cdot \frac{\mu_r \cdot \mu_0 \cdot h \cdot (r_2 - r_1)}{\pi \cdot (r_2 + r_1)}\\
		&= 250^2 \cdot \frac{500 \cdot 4\pi \cdot 10^{-7} \frac{\text{Vs}}{\text{Am}} \cdot 2 \cdot 10^{-2}\,\text{m} \cdot (5-3) \cdot 10^{-2}\,\text{m}}{\pi(5+3) \cdot 10^{-2}\,\text{m}}\\
		&= 62,5\,\text{mH}
		\end{align*}

		\item[b)]
		exakte Berechnung. Es wird berücksichtigt, dass die Flussdichte $B$ Radiusabhängig ist.\\
		1. Schritt: Feldstärke $H$
		\begin{align*}
		\oint\limits_s \vec{H}\cdot \d\vec{s} &= N \cdot I = H \cdot \ell \\
		H &= I \cdot \frac{N}{\ell} = I \cdot \frac{N}{2 \pi r}
		\end{align*}

		2. Schritt: Flussdichte $B$
		\begin{equation*}
		B = \mu_r \cdot \mu_0 \cdot H = \mu_r \cdot\mu_0 \cdot I \cdot \frac{N}{2 \pi r} \qquad\text{für } r_1 \ll r_2
		\end{equation*}

		3. Schritt: magnetischer Fluss $\Phi$
		\begin{align*}
		\Phi &= \int \vec{B} \cdot d\vec{A} \qquad\text{mit\,} d\vec{A} = h \cdot d\vec{r}\\
		&= \int\limits_{r_1}^{r_2} \vec{B} \cdot h \cdot d\vec{r} = \frac{\mu_r \cdot \mu_0 \cdot h \cdot N \cdot I}{2\pi} \cdot\int\limits_{r_1}^{r_2} \frac{1}{r} \cdot dr
		\end{align*}

		aus Grundintegral:
		\begin{equation*}
		\int\limits_{r_1}^{r_2} \frac{1}{r} \cdot dr =\ln(r_2)-\ln(r_1)= \ln \left(\frac{r_2}{r_1}\right)
		\end{equation*}

		folgt
		\begin{equation*}
		\Phi = \frac{\mu_r\cdot \mu_0 \cdot h \cdot N \cdot I}{2\pi} \cdot \ln \left(\frac{r_2}{r_1}\right)
		\end{equation*}

		4. Schritt: Induktivität $L$
		\begin{equation*}
		\text{Definition:\,}	L = \frac{N \cdot \Phi}{I}
		\end{equation*}

		Nach Einsetzen von $\Phi$ kürzt sich der Strom $I$ heraus.
		\begin{equation*}
		L_2 = N^2 \cdot \frac{\mu_r \cdot\mu_0 \cdot h}{2\pi} \cdot \ln \left(\frac{r_2}{r_1}\right)
		\end{equation*}

		Werte einsetzen:
		\begin{align*}
		L_2 &= 250^2 \cdot \frac{500 \cdot 4\pi \cdot 10^{-7} \frac{\text{Vs}}{\text{Am}} \cdot 2 \cdot 10^{-2} \,\text{m}}{2\pi} \cdot \ln \left(\frac{5\cdot 10^{-2}\,\text{m}}{3\cdot 10^{-2}\,\text{m}}\right)\\
		&= 63,85 \,\text{mH}
		\end{align*}

		\item[c)]
		Die Annahme eines homogenen Magnetfeldes im Kernquerschnitt trifft nur näherungsweise zu. Die Formel für $L_1$ ist deshalb nur eine gute Näherungsformel.
	\end{itemize}
}